% !TITLE 问刘十九
% !AUTHOR 白居易
% !DYNASTY 唐
% !GENRE 五言绝句
% !ORDER 48

\begin{poem}
绿蚁新醅酒\textsuperscript{1}，红泥小火炉。
晚来天欲雪，能饮一杯无\textsuperscript{2}？
\end{poem}

\begin{annotations}
\notetext{1}{绿蚁新醅酒：新酿好的酒上浮着微绿的泡沫，像小蚂蚁一样。醅（pēi），未过滤的酒。绿蚁，酒面上的绿色泡沫。}
\notetext{2}{无：文言疑问词，相当于“吗”。能不能来喝一杯呢？语气轻松如随口一问。}
\end{annotations}

\begin{appreciation}
《问刘十九》只有二十个字，是中国文学史上最著名的一次邀约之一。刘十九是白居易的朋友，生平不详——可有了这首诗，他的名字比许多宰相还长寿。醅是没过滤的酒，新酒面上浮着淡绿的泡沫，像小蚂蚁，所以叫绿蚁。

绿蚁新醅酒，红泥小火炉——新酿的酒，红泥的小火炉。一绿一红，一冷一暖，二十个字先铺两种颜色。酒是新的，火是刚烧旺的，一切都是正好”开始”的光景，就缺一个人。

晚来天欲雪，能饮一杯无——傍晚了，天要下雪了，你能不能来喝一杯？无就是”吗”：能喝一杯吗？全诗到此收住，一个问号悬在那里。刘十九来没来，白居易没说。可他越是轻描淡写，越让人心动——你想想，大雪天一个朋友问你”来不来喝一杯”，你不去才怪。

还记得《鹿鸣》吗？那是中国最早的宴饮诗，席上要鼓瑟吹笙、捧筐送礼，那是国宴。而《问刘十九》是宴饮的另一极：没有音乐，没有礼仪，没有第三个人，只有一炉火、一壶酒、一场雪。从《鹿鸣》到《问刘十九》，隔了一千七百年，中国人请客从请满堂宾客变成请一个知己，但”你来，我备了好东西”这份心意，一点没变。

白居易晚年住在洛阳，写了一大堆这样的小诗，不写大道理，只写小日子。有人嫌这些诗小，可”晚来天欲雪，能饮一杯无”被人念叨了一千多年——可见人心要的从来不是大场面，是有人惦记。

等你长大了，”能饮一杯无”翻译成你的话就是：下课了吗？要不要一起去吃碗热面？天冷了要记得约朋友，别总一个人闷着——朋友不是等来的，是你开口请来的。刘十九到底来没来，诗里没说；可要是换成你，我想你是会去的。
\end{appreciation}
